\section{Определение линейного пространства. Простейшие свойства. Примеры}

\begin{definition}
Множество $V$ называется \textbf{линейным (векторным) пространством} над числовым полем $F$ (где $F = \mathbb{R}$ или $\mathbb{C}$), если в $V$ введены две операции: сложение элементов и умножение элемента на число из $F$, удовлетворяющие следующим восьми аксиомам для любых $a, b, c \in V$ и $\lambda, \mu \in F$:
\begin{enumerate}
    \item $a + b = b + a$ (коммутативность сложения);
    \item $(a + b) + c = a + (b + c)$ (ассоциативность сложения);
    \item $\exists 0 \in V$ (нулевой вектор), такой что $\forall a \in V: a + 0 = a$;
    \item $\forall a \in V \; \exists (-a) \in V$ (противоположный вектор), такой что $a + (-a) = 0$;
    \item $1 \cdot a = a$;
    \item $\lambda(a + b) = \lambda a + \lambda b$ (дистрибутивность относительно сложения векторов);
    \item $(\lambda + \mu)a = \lambda a + \mu a$ (дистрибутивность относительно сложения чисел);
    \item $\lambda(\mu a) = (\lambda\mu)a$ (ассоциативность умножения на число).
\end{enumerate}
\end{definition}

\begin{theorem}[Простейшие свойства линейного пространства]
В любом линейном пространстве $V$ справедливы следующие свойства:
\begin{enumerate}
    \item Нулевой вектор $0$ и противоположный вектор $(-a)$ для любого $a$ определяются единственным образом.
    \item $0 \cdot a = 0$ для любого $a \in V$.
    \item $\lambda \cdot 0 = 0$ для любого $\lambda \in F$.
    \item $(-1) \cdot a = -a$ для любого $a \in V$.
\end{enumerate}
\end{theorem}

\begin{proof}
\begin{enumerate}
    \item Пусть существуют два нулевых вектора $0_1$ и $0_2$. Тогда $0_1 = 0_1 + 0_2 = 0_2 + 0_1 = 0_2$. Единственность противоположного вектора доказывается аналогично: если $a'$ и $a''$ --- противоположные к $a$, то $a' = a' + 0 = a' + (a + a'') = (a' + a) + a'' = 0 + a'' = a''$.
    \item Используя аксиомы 7 и 3: $0 \cdot a = (0 + 0) \cdot a = 0 \cdot a + 0 \cdot a$. Прибавив к обеим частям $-(0 \cdot a)$, получаем $0 = 0 \cdot a$.
    \item Используя аксиомы 6 и 4: $\lambda \cdot 0 = \lambda \cdot (a + (-a)) = \lambda a + \lambda(-a) = \lambda a + (-\lambda a) = 0$.
    \item Используя аксиомы 7, 5 и свойство 2: $(-1) \cdot a + a = (-1) \cdot a + 1 \cdot a = (-1 + 1) \cdot a = 0 \cdot a = 0$. Следовательно, $(-1) \cdot a$ является противоположным к $a$, то есть $(-1) \cdot a = -a$.
\end{enumerate}
\hfill$\blacksquare$
\end{proof}

\begin{example}[Примеры линейных пространств]
\begin{enumerate}
    \item Пространство $\mathbb{R}^n$ всех упорядоченных наборов из $n$ действительных чисел с покомпонентным сложением и умножением на число.
    \item Пространство $M_{m \times n}$ всех матриц размера $m \times n$ с элементами из $F$. Операции --- обычное сложение матриц и умножение матрицы на число.
    \[
    A = \begin{pmatrix}
    a_{11} & a_{12} & \dots & a_{1n} \\
    a_{21} & a_{22} & \dots & a_{2n} \\
    \vdots & \vdots & \ddots & \vdots \\
    a_{m1} & a_{m2} & \dots & a_{mn}
    \end{pmatrix} \in M_{m \times n}
    \]
\end{enumerate}
\end{example}
